\chapter{Algorithm and numerical safeguards}\label{ch:algorithm}

The solver, \code{complex\_spectrum}, isolates the zeros of an analytic
function $f$ in a rectangle $\Omega=[x_0,x_1]\times[y_0,y_1]$. In zero mode
$f=N$ (or the given function); in pole mode $f=D$. This chapter describes
each component and the conditions under which a result is declared
complete.

\section{Contour counts}\label{sec:count}

The boundary of a rectangle $B$ is sampled at $n$ equally spaced points per
edge, $4n$ points in total, traversed counterclockwise. With
$v_j=f(z_j)$, the phase and log-modulus increments are
\[
  \delta_j=\operatorname{Arg}\frac{v_{j+1}}{v_j}\in(-\pi,\pi],\qquad
  \mu_j=\log|v_{j+1}|-\log|v_j| ,
\]
and the discrete winding number is $w=\frac{1}{2\pi}\sum_j\delta_j$. The
number of samples starts at $n=n_0$ (\code{ContourPoints}, default 32) and
is doubled up to \code{ContourRefinements} times (default 7, i.e.\ up to
4096 points per edge). A count is \emph{accepted} when
\begin{enumerate}
  \item the rounded winding number $\operatorname{round}(w)$ is the same at
        three consecutive sampling levels,
  \item at the last two levels $|w-\operatorname{round}(w)|<10^{-7}$, and
  \item at the last two levels every increment satisfies $|\delta_j|<\pi/3$
        and $|\mu_j|<2$.
\end{enumerate}
Condition 3 is the practical form of the requirement that the argument of
$f$ vary slowly between samples; without it, a fast phase rotation could be
aliased and a whole turn lost. If any sample is zero, infinite or
\code{NaN} (a zero on the boundary, or overflow), or if the conditions are
not met at the finest level, the count fails and the rectangle is
\emph{unresolved}.

The same samples give the first-moment estimate~\eqref{eq:moment}, computed
as $\sum_j \tfrac{z_j+z_{j+1}}{2}(\mu_j+\ii\delta_j)/(2\pi\ii\,Z)$, used as a
Newton starting point.

\section{Newton's method}\label{sec:newton}

From a starting point $z$ inside a cell $B$ with count $m$, the solver
applies the modified Newton iteration
\[
  z\leftarrow z-\alpha\, m\,\frac{f(z)}{f'(z)} ,
\]
which converges quadratically to a zero of multiplicity $m$. Safeguards:
\begin{itemize}
  \item the step is limited to half the larger side of $B$;
  \item the damping factor $\alpha$ is halved (at most 12 times) until
        $|f|$ does not increase; if it cannot be decreased, the start is
        abandoned;
  \item the iteration stops if it leaves $B$, and converges when the step is
        below \code{RootTolerance}$\cdot(1+|z|)$.
\end{itemize}
The derivative is the user-supplied \code{Derivative}, the exact derivative
for symbolic input and coefficient vectors, or otherwise the four-point
complex difference
\[
  f'(z)\approx\frac{f(z+h)-f(z-h)-\ii\,[f(z+\ii h)-f(z-\ii h)]}{4h},
  \qquad h=10^{-4}\max(1,|z|),
\]
whose error is $O(h^4)$ for analytic $f$ (the $h^2$ terms of the real and
imaginary central differences cancel).

\section{Local confirmation}\label{sec:confirm}

A Newton limit $z^\ast$ is accepted only if a contour count on the small
square of half-width $\rho=\code{RootTolerance}\cdot(1+|z^\ast|)$ centered
at $z^\ast$ equals the count $m$ of the cell. This confirms both that
$z^\ast$ is a zero and that no other zero of the cell was missed; $m$ becomes
the multiplicity of $z^\ast$, and $\sqrt2\rho$ is reported as
\code{locationRadius}. Consequently, zeros closer than about $\rho$ are
reported as one location whose multiplicity is the size of the cluster.

\section{Recursive subdivision}\label{sec:subdivision}

\begin{algobox}{Algorithm 1 — recursive isolation of the zeros of $f$ in a cell $B$ with count $m$}
\begin{enumerate}[leftmargin=1.7em]
  \item If $m=0$, return.
  \item For each start in (first-moment center of $B$, geometric center of
        $B$, user seeds inside $B$): run Newton (Section~\ref{sec:newton});
        if the limit is confirmed with count $m$ (Section~\ref{sec:confirm}),
        record it and return.
  \item If the depth exceeds \code{MaxDepth} or the number of visited cells
        exceeds \code{MaxCells}, mark $B$ unresolved and return.
  \item Split $B$ across its longer side at the fraction $0.5$ of the side.
        Count the two halves. If both counts are accepted and add up to $m$,
        recurse into both halves and return.
  \item Otherwise retry with the split fractions $0.438\ldots$ and
        $0.573\ldots$, then with the other side. If no split conserves the
        count, mark $B$ unresolved.
\end{enumerate}
\end{algobox}

The \emph{conservation check} of step 4 detects zeros lying on the split
line (the counts of the halves then fail or do not add up) and moves the
line. The irrational-looking fractions make it unlikely that a second
attempt hits a zero again. Unresolved cells are never discarded: they are
returned in \code{info.unresolvedBoxes}.

\section{Cancellation}\label{sec:cancel}

In pole mode with a pair $N/D$, the zeros of $D$ are isolated with
Algorithm~1. For each location $z^\ast$ with multiplicity $m_D$, the zeros of
$N$ are counted on the same small square, giving $m_N$. The reported
multiplicity is $\max(m_D-m_N,0)$; if it is zero, $z^\ast$ is moved to
\code{cancelledLocations}, and the removed order $\min(m_D,m_N)$ is recorded.
Zero mode is symmetric. If the count of the other factor fails, the location
is kept and flagged in \code{cancellationUncertain}, and the result is not
complete. The test has the resolution of the small square: a pole and a zero
closer than $\rho$ are indistinguishable from an exact cancellation.

\section{Exploratory mode}\label{sec:exploratory}

When the function is not known to be analytic (an opaque handle without
\code{AssumeAnalytic}, or a single quotient handle), contour counts may
measure $Z-P$ instead of $Z$ and cannot be trusted. The solver then starts
Newton from a \code{GridSize} grid of points and from the seeds, and keeps
the limits confirmed by a positive local count. The result has status
\code{exploratory} and is never declared complete.

\section{The completeness contract}\label{sec:contract}

A search is \code{numerically\_complete} when (i) the count on
$\partial\Omega$ was accepted, (ii) every cell was resolved, (iii) the
multiplicities of the located zeros add up to that count, and (iv) every
cancellation test was resolved. Under the analyticity assumption, and if
the accepted counts are correct, all zeros in $\Omega$ are then reported,
each within its \code{locationRadius} as a cluster of the stated
multiplicity.

\begin{keybox}{What this is not}
It is not a proof. The counts are computed from finitely many floating-point
samples; a function that oscillates between samples faster than the
safeguards of Section~\ref{sec:count} can detect could in principle fool
them. Analyticity of a handle is the user's declaration. Residuals depend on
the scaling of the searched factor and are not bounds on the location
error. \code{info.certified} is always false. Rigorous enclosures would
require interval arithmetic, which is outside the scope of this version.
\end{keybox}

\section{Cost}

The cost is dominated by function evaluations on contours: $4n$ per count,
with $n$ doubling until acceptance, for the outer rectangle, for every
cell of the subdivision and for every local confirmation. Vectorized
handles are evaluated on all samples of a contour at once. Tall rectangles
over oscillatory functions (e.g.\ $e^{-sT}$ with large $T|\Imag s|$) need
more samples per edge; smaller rectangles or larger
\code{ContourRefinements} help. Seeds from a previous parameter value
reduce the number of subdivisions in parameter sweeps.

\section{Related methods and software}\label{sec:related}

Contour-integral root finders go back to \citet{delves1967}; see
\citet{kravanja2000} for a thorough treatment including moment-based
eigenvalue formulations. \code{cxroots} \citep{cxroots} implements such a
method in Python with subdivision. GRPF \citep{grpf} analyzes the phase of a
single meromorphic function on an adaptive triangulation and locates both
zeros and poles. For quasi-polynomials and time-delay systems, QPmR
\citep{vyhlidal2009} maps the zero-level curves of the real and imaginary
parts, and spectral discretization methods \citep{breda2005,breda2015,michiels2014}
approximate the rightmost roots of matrix delay equations. ContourRoots
combines well-known ingredients; its emphasis is on a MATLAB-style
interface, explicit cancellation for $N/D$ pairs, and diagnostics that make
incomplete searches visible.
